%%
%% This is file `main.tex' based on `sample-sigconf.tex' (q.v. for spurce of that,
%%
%% IMPORTANT NOTICE:
%% 
%% For the copyright see the original source file `sample-sigconf.tex'
%% in the `Sample' folder.
%%
%% For distribution of the original source see the terms
%% for copying and modification in the file samples.dtx.
%% 
%% This generated file may be distributed as long as the
%% original source files, as listed above, are part of the
%% same distribution. (The sources need not necessarily be
%% in the same archive or directory.)
%%
%% Commands for TeXCount
%TC:macro \cite [option:text,text]
%TC:macro \citep [option:text,text]
%TC:macro \citet [option:text,text]
%TC:envir table 0 1
%TC:envir table* 0 1
%TC:envir tabular [ignore] word
%TC:envir displaymath 0 word
%TC:envir math 0 word
%TC:envir comment 0 0
%%
%%
%% The first command in your LaTeX source must be the \documentclass command.

%% NOTE that a single column version is required for 
%% submission and peer review. This can be done by changing
%% the \doucmentclass[...]{acmart} in this template to 
%%\documentclass[manuscript,review,anonymous]{acmart}
%% This version is used for drafting and final submission
\documentclass[sigconf]{acmart}


%% 
%% To ensure 100% compatibility, please check the white list of
%% approved LaTeX packages to be used with the Master Article Template at
%% https://www.acm.org/publications/taps/whitelist-of-latex-packages 
%% before creating your document. The white list page provides 
%% information on how to submit additional LaTeX packages for 
%% review and adoption.
%% Fonts used in the template cannot be substituted; margin 
%% adjustments are not allowed.

%%
%% \BibTeX command to typeset BibTeX logo in the docs
\AtBeginDocument{%
  \providecommand\BibTeX{{%
    \normalfont B\kern-0.5em{\scshape i\kern-0.25em b}\kern-0.8em\TeX}}}

%% Rights management information.  This information is sent to you
%% when you complete the rights form.  These commands have SAMPLE
%% values in them; it is your responsibility as an author to replace
%% the commands and values with those provided to you when you
%% complete the rights form.
\setcopyright{acmlicensed}
\copyrightyear{2024}
\acmYear{2024}
\acmDOI{XXXXXXX.XXXXXXX}

%% These commands are for a PROCEEDINGS abstract or paper.
\acmConference[Conference acronym 'XX]{Make sure to enter the correct
  conference title from your rights confirmation email}{June 03--05,
  2018}{Woodstock, NY}
%
%  Uncomment \acmBooktitle if th title of the proceedings is different
%  from ``Proceedings of ...''!
%
%\acmBooktitle{Woodstock '18: ACM Symposium on Neural Gaze Detection,
%  June 03--05, 2018, Woodstock, NY} 
\acmISBN{978-1-4503-XXXX-X/18/06}


%%
%% Submission ID.
%% Use this when submitting an article to a sponsored event. You'll
%% receive a unique submission ID from the organizers
%% of the event, and this ID should be used as the parameter to this command.
%%\acmSubmissionID{123-A56-BU3}

%%
%% For managing citations, it is recommended to use bibliography
%% files in BibTeX format.
%%
%% You can then either use BibTeX with the ACM-Reference-Format style,
%% or BibLaTeX with the acmnumeric or acmauthoryear sytles, that include
%% support for advanced citation of software artefact from the
%% biblatex-software package, also separately available on CTAN.
%%
%% Look at the sample-*-biblatex.tex files for templates showcasing
%% the biblatex styles.
%%

%%
%% The majority of ACM publications use numbered citations and
%% references.  The command \citestyle{authoryear} switches to the
%% "author year" style.
%%
%% If you are preparing content for an event
%% sponsored by ACM SIGGRAPH, you must use the "author year" style of
%% citations and references.
%% Uncommenting
%% the next command will enable that style.
%%\citestyle{acmauthoryear}

%%
%% end of the preamble, start of the body of the document source.

%% % Location of your graphics files for figures, here a sub-folder to the main project folder
\graphicspath{{./images/}} 


\begin{document}

%%
%% The "title" command has an optional parameter,
%% allowing the author to define a "short title" to be used in page headers.
\title{Hierarchical Deep Learning for RNA Motif Classification from Cryo-EM Density Maps}

%%
%% The "author" command and its associated commands are used to define
%% the authors and their affiliations.
%% Of note is the shared affiliation of the first two authors, and the
%% "authornote" and "authornotemark" commands
%% used to denote shared contribution to the research.
\author{Ketsia Mbaku}
\email{ketsia@uw.edu}
\affiliation{%
  \institution{University of Washington, Bothell}
  \streetaddress{18115 Campus Way NE}
  \city{Bothell}
  \state{Washington}
  \country{USA}
  \postcode{98011}
}

%%
%% By default, the full list of authors will be used in the page
%% headers. Often, this list is too long, and will overlap
%% other information printed in the page headers. This command allows
%% the author to define a more concise list
%% of authors' names for this purpose.
\renewcommand{\shortauthors}{Mbaku}

%%
%% The abstract is a short summary of the work to be presented in the
%% article.
\begin{abstract}
Accurate classification of RNA secondary structural motifs from
cryo-EM density maps is a key step toward automated RNA structure
prediction pipelines. Prior work established that broad motif
categories can be predicted from three-dimensional density volumes,
but fine-grained classification across the full 25-class CoSSMos
taxonomy has remained out of reach, with flat supervision yielding
only 38.4\% accuracy in our preliminary baseline experiments. We
attribute this failure to the absence of structural guidance in
the training objective: a flat cross-entropy loss treats all
confusions equally, offering no signal that reflects the biological
organization of the label space. In this work we propose a hierarchical classification
framework in which motifs are first distinguished by topology,
then by strand symmetry, and finally by size, mirroring the
natural structure of the CoSSMos taxonomy. We evaluate three deep
learning models of increasing sophistication on a resolution-filtered
subset of a large-scale cryo-EM RNA motif dataset: a flat 3D U-Net
encoder retrained as a 25-class baseline, a 3D ResNet-18 trained
with a multi-head hierarchical loss, and a 3D Swin Transformer
with conditional hierarchical inference in which each prediction
level gates the next. Across all 25 classes, an average of 87.4\%
of dataset instances satisfy our 5.0~\AA{} resolution criterion,
with class-weighted training applied to compensate for the
underrepresented classes that fall below this average. We expect
hierarchical supervision to improve
accuracy and to reduce the within-family confusion
errors that dominate the baseline, with the
Swin Transformer providing additional gains through
spatial reasoning over the density volume.
\end{abstract}

%%
%% The code below is generated by the tool at: http://dl.acm.org/ccs.cfm
%% Please copy and paste the code instead of the example below.
%%


%% The following are not a requirement, delete if not using
\received{12 May 2026}  %% inital submission date

%%
%% This command processes the author and affiliation and title
%% information and builds the first part of the formatted document.
\maketitle
\section{Introduction}
% ============================================================
 
\subsection{Background and Motivation}
 
Ribonucleic acid plays a central role in a broad range of cellular
processes, from catalyzing biochemical reactions and regulating gene
expression to serving as the primary genetic material in many viruses
\cite{tinoco1999}. Unlike the rigid double-helical structure
of DNA, single-stranded RNA folds back on itself through intramolecular
base pairing, producing a rich landscape of three-dimensional
conformations. These conformations are not random: the same local
structural patterns, known as secondary structural motifs, recur
across evolutionarily distant RNA molecules and are tightly coupled
to biological function \cite{leontis2001}. Because many of these
motifs serve as binding sites for small molecules and therapeutic
agents, their accurate identification and classification is of direct
relevance to structure-based drug discovery \cite{porter2017}.

A key limitation of current approaches is that they operate at the
level of full reconstruction, attempting to trace backbone
coordinates across an entire map before any motif-level understanding
is established. This top-down strategy is particularly fragile for RNA,
where flexible single-stranded regions and repetitive structural elements
make global tracing error-prone. An alternative strategy, motivated by
the modular nature of RNA architecture, is to first classify and localize
recurring local motifs from density data, and then use these intermediate
labels to guide downstream structure prediction. Realizing this strategy
requires machine learning models capable of recognizing
structural categories directly from three-dimensional density volumes.
\subsection{Related Work}
A critical prerequisite for motif classification from
cryo-EM is the availability of a dataset pairing density with motif-level structural labels.
\citeauthor{murugadass2026} \cite{murugadass2026} recently introduced
such a resource: a dataset of over 125,000 motif-resolved cryo-EM
density maps spanning the full 25-class CoSSMos taxonomy, covering
resolutions from 1.5 to 34~\AA{}. Each instance is represented as a
standardized 64$\times$64$\times$64 voxel patch centered on the motif,
with voxel-level labels for backbone, ribose sugar, and nucleobase
components derived from aligned atomic models. To demonstrate the
dataset's utility for supervised learning, the authors trained a 3D convolutional neural network comprising three
convolutional blocks on a subset of high-resolution maps
(1.5--2.8~\AA{}), collapsing the 25 fine-grained CoSSMos classes into
five groups: symmetric internal loops, asymmetric internal loops,
hairpins, bulges, and a background class. At this coarse granularity,
the classifier achieved a macro-average specificity of 0.948, confirming
that the density signal is sufficient to discriminate broad structural
categories. However, the transition from 25 fine-grained to 5 coarse
classes did not explicitly encode the
symmetric-to-asymmetric distinction within internal loops, and
discards the size-level granularity that distinguishes motifs within
each family. Whether fine-grained classification across the full
CoSSMos taxonomy is achievable from density data alone was left as
an open question by the authors.
 
In a prior work as part of CSS581 Machine Learning course during the Fall 2025, we conducted a comparison of classical
and deep learning classifiers for RNA motif classification from
cryo-EM density, using a subset of 15 symmetric motif classes (five
hairpin lengths, five bulge sizes, and five symmetric internal loop
sizes) drawn from the same dataset. A CNN model was evaluated at both
6-class and 15-class granularity. At coarse granularity,
the model exceeded 92\% accuracy, confirming the feasibility of
broad category prediction. However, at fine-grained granularity the
picture diverged sharply: the 3D CNN collapsed to 55.69\% accuracy. Confusion analysis revealed that errors were
concentrated within motif families, primarily between size-adjacent
classes such as bulge3 and bulge4. Critically, the
flat 15-class label space provided no structural guidance to the
model: a confusion between bulge2 and bulge3 was penalized identically
to a confusion between a bulge and a hairpin, despite the former being
a far less serious error from a biological standpoint.
 
This literature review reveals three gaps that motivate
the present work. First, no prior study has demonstrated
fine-grained classification across the full 25-class CoSSMos taxonomy
from cryo-EM density. The closest attempts either stopped at 5
coarse classes \cite{murugadass2026} or at 15 symmetric classes
\cite{mbaku2025}. Second, both prior works
employed label spaces that ignore the biological hierarchy
inherent in the CoSSMos taxonomy. Third, more powerful  architectures
capable of capturing spatial dependencies within density
patches, such as residual networks and vision transformers, have not
been evaluated in this domain.

\subsection{Research Questions and Hypotheses}
This work evaluates whether training a model to predict motif labels at multiple levels of biological detail, rather than all 25 classes at once, leads to better fine-grained classification. We hypothesize that learning broad categories first and progressively refining toward size-level distinctions will give the model structural guidance that a flat training objective cannot provide.
We also ask whether this hierarchical approach specifically fixes the confusion patterns observed in prior work, particularly the mixing of symmetric and asymmetric internal loops and errors between size-adjacent classes within the same family. We hypothesize that supervising the symmetry level explicitly will force the model to resolve that distinction early, before attempting the harder size-level predictions.
Finally, we aim to explore whether making predictions conditionally, where each level informs the next, works better than predicting all levels independently in parallel, and whether a transformer-based architecture captures features that convolutional models miss. We expect both to contribute, especially on the rarest and most structurally similar classes.


% ============================================================
\section{Methods}
% ============================================================
 
\subsection{Data Source and Scope}

This work builds on the large-scale cryo-EM RNA motif dataset
introduced by \citeauthor{murugadass2026} \cite{murugadass2026},
which provides over 125,000 motif-resolved density map segments
spanning all 25 RNA CoSSMos \cite{cossmosdb} secondary structural motif classes \cite{murugadass2026dataset}. Each segment is a standardized
$64 \times 64 \times 64$ voxel patch extracted from experimental
EMDB maps \cite{emdb2023} and paired with its corresponding atomic
model from the RCSB Protein Data Bank \cite{rcsb2019}. The full
dataset covers a resolution range of 1.5 to 34.0~\AA{}, but
structural detail degrades substantially beyond 5~\AA{}, with
cross-correlation scores dropping from the 0.70--0.75 range at
3.0--4.0~\AA{} to the 0.50--0.60 range at 5.0--10.0~\AA{}
\cite{murugadass2026}. To ensure that density features are
sufficiently resolved to support fine-grained classification, we
restrict our working dataset to maps with a reported global
resolution of 5.0~\AA{} or better.

To characterize the usable dataset under this constraint, we
developed two preprocessing scripts. The first queries the metadata
of each segmented motif instance to retrieve the reported resolution
of its source EMDB map, retaining only instances at or below
5.0~\AA{}. The second tabulates per-class instance counts within
the filtered set to identify any classes at risk of severe
underrepresentation. Figure~\ref{fig:resolution_histogram} shows
the percentage of maps meeting the resolution criterion for each
of the 25 motif classes. Across all classes, an average of 87.4\%
of instances satisfy the threshold, indicating that the resolution
filter preserves the large majority of the dataset. Most classes
retain above 85\% of their instances, with several exceeding 95\%,
including the 1$\times$5 asymmetric loop (96.7\%), bulge5 (97.7\%),
and bulge4 (96.0\%). However, three classes fall below the overall
average by a notable margin: bulge1 retains only 66.2\% of its
instances, hairpin4 retains 66.3\%, and hairpin7 retains 71.6\%.
The 3$\times$5 and 4$\times$4 asymmetric loop classes also show
modest shortfalls at 83.7\% and 79.6\% respectively. To mitigate
the impact of these gaps, class-weighted loss and stratified
sampling will be applied during training, ensuring that
underrepresented classes receive proportionally greater emphasis
without being discarded from the evaluation.

\begin{figure}[h]
    \centering
    \includegraphics[width=\columnwidth]{histogram.png}
    \caption{Percentage of motif instances per class satisfying the
    5.0~\AA{} resolution threshold across all 25 CoSSMos motif
    classes.}
    \label{fig:resolution_histogram}
\end{figure}
\subsection{Baseline Model}
\label{sec:baseline}
% ============================================================
 
The baseline model is a direct extension of the architecture
developed in prior CSS581 work, retrained here on the
full 25-class label space to establish the performance ceiling
achievable without hierarchical supervision. The architecture
is a CNN with three downsampling stages, processing
a normalized $64 \times 64 \times 64$ density patch as a
single-channel input. Each downsampling stage applies a pair of
$3 \times 3 \times 3$ convolutions with batch normalization and
ReLU activations, followed by $2 \times 2 \times 2$ max pooling.
Channel dimensions progress from 1 to 32, 64, 128, and 256,
with adaptive average pooling applied at the bottleneck to produce
a fixed 256-dimensional density embedding independent of voxel
spacing. In the original work \cite{mbaku2025}, this density
embedding was fused with hand-crafted sequence and base-pairing 
features obtained from PDB via parallel MLP branches before classification. In this
work, we remove those feature branches and train the
model on density alone, for consistency with the subsequent models to be trained and
to isolate the contribution of the density across all
three architectures.
 
Preliminary results from retraining this architecture on the
25-class flat objective are summarized in Table~\ref{tab:baseline_results}.
As anticipated, extending from the 15-class setting of prior work
to the full 25-class taxonomy produces a marked degradation in
classification performance. Overall test accuracy falls to 38.4\%,
compared to 55.69\% on 15 classes, reflecting
the additional difficulty introduced by the 10 asymmetric internal
loop subclasses, which share highly similar density signatures
and are severely underrepresented relative to the bulge and hairpin
families. Macro-average sensitivity drops to 0.31, indicating that
the model systematically fails to recover many fine-grained classes. These results
confirm that flat supervision at the 25-class level is insufficient
and motivate the hierarchical training strategies to be evaluated in
subsequent models as part of this project.
 
\begin{table}[h]
\caption{Baseline model performance on the flat 25-class
classification task}
\label{tab:baseline_results}
\centering
\begin{tabular}{lc}
\toprule
Metric & Value \\
\midrule
Overall accuracy         & 38.4\% \\
Macro sensitivity        & 0.31 \\
Macro specificity        & 0.94 \\
Macro AUC (one-vs-rest)  & 0.84 \\
\midrule
Sensitivity -- hairpin3  & 0.74 \\
Sensitivity -- hairpin7  & 0.61 \\
Sensitivity -- bulge2    & 0.78 \\
Sensitivity -- bulge5    & 0.43 \\
Sensitivity -- 1$\times$1 symmetric & 0.59 \\
Sensitivity -- 5$\times$5 symmetric & 0.34 \\
Sensitivity -- 1$\times$2 asymmetric & 0.27 \\
Sensitivity -- 4$\times$5 asymmetric & 0.08 \\
\bottomrule
\end{tabular}
\end{table}
 
\subsection{Proposed Hierarchy and Project Plan}
 
We organize the 25 CoSSMos classes into a three-level biological
hierarchy that forms the backbone of both the training objectives
and the experimental progression. At the first level (L1), motifs
are grouped by topology into hairpins, internal loops, and bulges,
reflecting the most geometrically distinct density signatures. At
the second level (L2), internal loops are further split by strand
symmetry into symmetric and asymmetric subclasses, a distinction
never explicitly encoded in prior work and the locus of the hardest
confusions in \cite{mbaku2025}. The third level (L3) retains all
25 fine-grained CoSSMos classes and is the primary evaluation
target; L1 and L2 serve exclusively as intermediate supervision
signals. The full mapping is shown in Table~\ref{tab:hierarchy}.
 
\begin{table}[h]
\caption{Three-level hierarchy over all 25 CoSSMos motif classes.}
\label{tab:hierarchy}
\centering
\begin{tabular}{lll}
\toprule
L1 (Topology) & L2 (Symmetry) & L3 (25 classes) \\
\midrule
Hairpin       & Hairpin        & hairpin3--hairpin7 \\
Internal loop & Symmetric      & 1$\times$1 -- 5$\times$5 \\
Internal loop & Asymmetric     & 1$\times$2, \ldots, 4$\times$5 \\
Bulge         & Bulge          & bulge1--bulge5 \\
\bottomrule
\end{tabular}
\end{table}
 
The remaining work proceeds in four tasks. The first is to complete
the preprocessing pipeline: resolution filtering and per-class
counting scripts have been developed and validated; full execution
across the corpus is ongoing, after which stratified splits will
be locked and shared across all models. The second task is to
implement a 3D ResNet-18 with three parallel
classification heads trained jointly under the weighted hierarchical
loss $\mathcal{L} = \alpha\mathcal{L}_{L1} + \beta\mathcal{L}_{L2}
+ \gamma\mathcal{L}_{L3}$, with a grid search over the weighting
hyperparameters. The third task is a 3D Swin Transformer
\cite{swin2021} with conditional hierarchical inference, where the
L1 softmax output gates the L2 head and the L2 output gates L3,
enforcing a decision path consistent with the biological hierarchy;
gradient checkpointing will be applied if memory constraints
require it. The fourth task is a unified comparative evaluation
on the held-out test set, including per-class sensitivity across
all 25 classes, confusion matrix analysis at each hierarchy
level, and an ablation study isolating the contribution of the
hierarchical loss from the architectural improvement in the second model.

\section{Conclusion}

This report presents a progress update toward hierarchical deep
learning classification of RNA secondary structural motifs from
cryo-EM density maps. Building on prior work that identified the
failure of flat supervision at fine-grained granularity, we
introduced a three-level biological hierarchy over all 25 CoSSMos
motif classes and proposed three models designed to exploit it
progressively. Preliminary baseline results confirm that flat
25-class training is insufficient, achieving only 38.4\% accuracy
with systematic failure on the asymmetric internal loop family,
establishing a clear performance floor that the hierarchical
models are designed to surpass. Data preprocessing results show
that 87.4\% of available instances satisfy the 5.0~\AA{} resolution
criterion, providing a representative
training set across all motif classes.

The remaining work centers on implementing and evaluating the
3D ResNet-18 with hierarchical loss and the 3D Swin Transformer
with conditional inference. Together these models will test whether
hierarchical supervision recovers fine-grained accuracy, whether
it specifically repairs the symmetry-level confusions documented
in the baseline, and whether attention-based spatial reasoning
provides complementary gains over convolutional architectures.
Beyond the immediate experimental goals, this work constitutes
a
benchmark that can serve as a foundation for future motif-aware
RNA structure prediction methods.
\bibliographystyle{ACM-Reference-Format}
\bibliography{references.bib}
\end{document}
\endinput
%%
%% End of file `main.tex'.
